\begin{frame}{Pendulum 학습 결과 (300 반복 $\times$ 400 스텝)}
  \begin{columns}[T]
    \begin{column}{0.56\textwidth}
      \centering
      \includegraphics[width=\linewidth]{ch11_2_ppo_pendulum_curve.png}
      {\footnotesize 시드 42, 총 12만 스텝. 연회색 = 에피소드별
      리턴(항상 음수, 0에 가까울수록 좋음), 파랑 실선 = 25에피소드
      이동평균.}
    \end{column}
    \begin{column}{0.42\textwidth}
      {\small
      \begin{tabular}{@{}lr@{}}
        \toprule
        처음 10 에피소드 평균 & $-805.3$ \\
        마지막 10 에피소드 평균 & $-742.3$ \\
        \bottomrule
      \end{tabular}
      \vskip0.5em
      \begin{itemize}
        \item 약 \textbf{8\%} 개선, 여전히 절댓값이 크고 학습
          곡선도 매끄럽지 않음 (중간 지점에서 오히려 나빠지는
          구간도 있었음).
        \item 구현 오류가 아니라, Pendulum 같은 연속 제어는
          완전 수렴에 \textbf{수십만$\sim$수백만 스텝}이 필요
          — 정직한 결과.
        \item DQN(CartPole, 300에피소드 만에 상당한 개선, 9.3절)
          과 비교: 연속 제어가 이산보다 \textbf{더 오래} 걸림.
      \end{itemize}}
    \end{column}
  \end{columns}
\end{frame}
